
\chapter{Part II Review: Actions, Waves, and Conservation}

The action principle turns a global quantity into a local equation.  The required step is always the same: vary the object, integrate by parts, state the boundary condition, and use the arbitrariness of the remaining variation.

\section{A guided reconstruction}

% QFTFIGURE-BEGIN partreview02-01-part-map
\qftfigure{figures/instructional/partreview02/partreview02-01-part-map.pdf}{The map collects the part's main constructions into a visual route so the reader can see how the chapters support one another.}{fig:partreview02-01-part-map}
% QFTFIGURE-END partreview02-01-part-map












For a path \(q(t)\), stationarity gives the Euler-Lagrange equation.  For a field \(\phi(t,\bm x)\), the same derivation has one integration by parts for each coordinate:
\[
  \frac{\p\Lag}{\p\phi}-\p_\mu\frac{\p\Lag}{\p(\p_\mu\phi)}=0.
\]
When the action is invariant under a continuous transformation, the coefficient of the local version of the transformation is a conserved current.

\begin{exercise}
Derive the field Euler-Lagrange equation from an action \(S=\int\dd^4x\,\Lag(\phi,\p_\mu\phi)\).
\begin{answer}
Vary \(\phi\), integrate the \(\p_\mu\delta\phi\) term by parts, discard the stated surface term, and set the coefficient of arbitrary \(\delta\phi\) to zero.
\end{answer}
\end{exercise}
